\scanpage{20}\mspage{18}
\deleted{\unsure{donc}} \unsure{i.e.}
\[
\operatorname{card}E_\varphi=2(\deleted{\unsure{l}}k-\deleted{\unsure{k}}l)
\]
Il est immédiat qu'un foncteur $\varphi$ est surjectif
sur les objets et sur les flèches, donc identifie
$\unsure{\mathscr{L}|\mathcal{L}}_l$ à une catégorie quotient de $\mathscr{L}_k$ -- \deleted{\unsure{savoir}}
et la relation d'équivalence dans $\mathscr{L}_k$ définie
par $\varphi$ est celle obtenue en \deleted{\unsure{écrivant que les}} \inserted{identifiant}
tous les objets qui sont sources ou buts de
flèches d'une $=$ \unsure{composante|compos.te} connexe \inserted{$I$} de $E_\varphi$
\diagram{\begin{tikzcd}[ampersand replacement=\&,column sep=small]
\cdots \& {} \arrow[r,phantom] \& {} \arrow[l] \arrow[r] \& \bullet \& \bullet \arrow[l] \arrow[r] \& \bullet \& {} \arrow[l] \arrow[r] \& {} \& \cdots
\arrow[from=1-3,to=1-7,phantom,"{\underbrace{\hspace{6em}}_{I}}"']
\end{tikzcd}}
et les flèches dans cette composante à l'identité correspondante. Les conditions énoncées sur $E$
(indépendamment d'une $\varphi$ qui lui \unsure{donnerait}
\unsure{nai ssance|naissance}) assurent que la catégorie quotient
est isomorphe \inserted{(avec \uline{préservation} de source et de but)}
soit à $\mathscr{L}_l$, soit à la duale $\unsure{\breve{\mathscr{L}}|\check{\mathscr{L}}}_l=\mathscr{L}_l^{\unsure{\vee|\smile}}$,
suivant la direction de la première flèche de
\deleted{c'est le premier si $k$ et $l$ sont de même parité, et
le deuxième \deleted{\unsure{dans le}} cas contraire. (Car en
tous les cas, la catégorie quotient ``commence'' par
une flèche de $=$ direction que $\mathscr{L}_k$, laquelle
en dépend \unsure{pour} de la parité)}
$\overline{\mathscr{L}}_k$ (quotient de $\mathscr{L}_k$).
Cette ``direction \inserted{initiale}'' est la $=$ que celle pour $\mathscr{L}_k$ si
la première flèche \inserted{de $\mathscr{L}_k$} n'appartient pas à $E$, ou bien
si \deleted{sa comp.} elle y appartient et sa comp. conn.
dans $E$ est de cardinal pair. Dans le cas
inverse ($u_{\unsure{1|*}}\in E$ et sa comp. conn. dans $E$ est de
cardinal impair), \deleted{elle est de cont.} \inserted{la direction initiale de $\overline{\mathscr{L}}_k$ est opposée à celle de $\mathscr{L}_k$.}
% archive sheet number: 19
